\documentclass[12pt,a4paper]{article} %12 кегль, формат бумаги А4, тип статья

\usepackage[T2A]{fontenc} %поддержка кириллицы
\usepackage[utf8]{inputenc} %кодировка символов юникод
\usepackage[english,russian]{babel} %список рабочих языков
\usepackage{cmap} %для поиска в pdf
\usepackage[left=2cm,right=2cm,top=2cm,bottom=2cm]{geometry} %размеры отступов от краев страницы
\usepackage{graphicx} %возможность работать с графикой
\usepackage{enumitem} %для работы со списками
\usepackage{float} %для точного размещения таблиц и графики
\usepackage{caption} %для создания подписей таблиц и графики
\usepackage{amsmath,amssymb,amsthm,amsfonts} %штуки для математики
\usepackage{extarrows} %чтобы подписи над и под нестандартными стрелками писать короче в коде
\usepackage{makeidx} %собирает ключевые слова для предметного указателя
\usepackage[hidelinks]{hyperref} %убирает выделение ссылок; добавляет гиперссылки в оглавление, перекрестные ссылки и URL
\usepackage{setspace} %устанавливает отступы
\linespread{1.5} %длина межстрочного интервала
\vspace{125mm} %длина абзацного отступа
\usepackage{titlesec}
\titleformat{\section}{\normalfont\Large\bfseries\centering}{}{0pt}{} % эти две строчки убирают циферки перед названиями разделов в основном тексте и выравнивают названия секций по центру
\usepackage{paratype} %чтобы точно был поиск и шрифт хороший

\author{Автор конспекта: Краснов Борис} %автор
\title{Жеданов А.С.: Математическая физика} %название
\date{Осень 2026} %дата создания документа

\newtheorem{theorem}{Теорема}
\newtheorem{lemma}{Лемма}
\newtheorem{definition}{Определение}
\newtheorem{proposition}{Предложение}
\newtheorem{consequence}{Следствие}
\newcommand{\R}{\mathbb{R}}
\newcommand{\Co}{\mathbb{C}}
\newcommand{\E}{\mathbb{E}}
\newcommand{\D}{\mathbb{D}}
\newcommand{\PR}{\mathbb{P}}
\newcommand{\charf}[2][t]{\varphi_{#2}(#1)}
\newcommand{\F}{\mathcal{F}}
\newcommand{\M}{\mathfrak{M}}

\begin{document}
	\maketitle
	\tableofcontents
	\clearpage
	
	\section{Слишком долгое введение}
	Вообще матфиз в простом понимании - всё, ну или если культурнее, то этот раздел математики охватывает большинство ее частей. У нас же в основном будут изучаться дифференциальные уравнения в частных производных. Также мы будем большое внимание уделять физическим аспектам, ибо мы математики и физики у нас мало, поэтому вот, будет больше.
	
	Рассмотрим некоторые примеры задач. Первый - всем известное уравнение Лапласа:
	$$\triangle F(z)=0,$$
	где уравнение рассматривается в некоторой области из $\R^n$, а $\triangle$ - оператор Лапласа.
	
	Можно, вспоминая лекцию по вариационке, дать определение лапласиану, так как это сделал Роман Владимирович, а можно по-человечески:
	$$\triangle=\sum_{k=1}^n\frac{\partial^2}{\partial x_k^2}$$
	
	Это пример уравнения эллиптического типа.
	
	Физический смысл данного уравнения: есть в пространстве электрическое поле $\overrightarrow{E}$ и в каждой точке есть потенциал поля - некоторая функция:
	$$\nabla\varphi=\overrightarrow{E}$$
	Как вычислять изменение потенциала мы знаем - интеграл по кривой между точками:
	$$\int\limits_{\gamma}\overrightarrow{E}ds=\varphi(B)-\varphi(A)$$
	Прикол данного поля - изменение не зависит от кривой. Такие поля называются консервативными или потенциальными.
	
	Отсюда можно определить потенциаль в точке, положив по определению потенциал в бесконечности 0:
	$$\int_A^{\infty}\overrightarrow{E}ds=-\varphi(A)$$
	
	Ну и естественная задача такая: пусть в некоторой области пространства известно распределение заряда, а мы хотим найти значение потенциала в пространстве. Тогда надо решить уравнение Лапласа в данной области.
	
	Теперь рассмотрим уравнение Гельмгольца:
	$$\triangle F(z)=\kappa F(z),$$
	где $\kappa$ - некоторый фикс. коэффициент.
	
	Его физический смысл - описание распределения волн в среде, если они гармонически зависят от времени.
	
	Также есть волновое уравнение, относящееся к типу гиперболических:
	$$\square F(z,t=0),$$
	где $\square:=\triangle-\frac{1}{c^2}\partial_t^2$ - оператор Даломбера (D'Alembert), а $c$ некоторая константа, но может зависить от точки пространства. Если посмотреть на применение этого уравнения в физики, то можно догадаться, что эта константа отвечает скорости, ибо иначе что-то странное с размерностями.
	
	Ну и последний тип уравнений - параболические. Например уравнение теплопроводности:
	$$\partial_t F(z,t)=b\triangle F(z,t)$$
	Физический смылс понятен - брусочки разной температуры приложили друг к другу и надо просмотреть изменение температуры.
	
	Ну и немножко квант. меха - уравнение Шредингера:
	$$i\partial_t\psi=(E-U)\triangle\psi,\quad U=U(z)$$
	Отличие этого уравнения от предыдущего - комплексность, а если посмотреть на физику, то система описываемая предыдущим придет в равновесие и будет иметь промежуточное состояние от исходных, а в квант. мехе $\psi$ описывает некоторое свойство некой характеристики среды и в процессе система не стабилизируется, а будет иметь колебательный характер.
	
	Есть у нас еще и электро-магнитное поле со своими уравнениями. Соответственно есть электрическое поле $\overrightarrow{E}$ и магнитное $\overrightarrow{H}$ Вот некоторые уравнения:
	$$div\overrightarrow{E}=4\pi\rho$$
	$$div\overrightarrow{H}=0$$
	Дивергенция описывает поток выходящего вещества из некоторой области и да, это вектор, ну так, на всякий.
	
	Далее был разговор про гильбертовы пространства и мы все тут знаем, так что напишу лишь то, что вообще знать не нужно, и соответственно не все знают. Если мы рассмотрим отрезок $[-1,1]$ и ортонормируем там систему $\{1,x,x^2,\cdots\}$, то полученный базис будет состоять из полиномов, которые называются полиномами Лежандра.
	\section{Обобщенные функции}
\end{document}